\documentclass[tikz,border=4pt]{standalone}
\usepackage{ctex}
\usepackage{amsmath}
\usepackage{pgfplots}
\pgfplotsset{compat=1.18}
\usetikzlibrary{calc, arrows.meta}

\begin{document}
% part14/03 等腰直角三角形3种情形
% 抛物线 y=x^2-2x-3, A(-1,0), B(3,0)
% 情形1:直角在A (P在x=-1竖线上); 情形2:直角在B (P在x=3竖线上); 情形3:直角在P
\begin{tikzpicture}
\begin{axis}[
  axis lines=center,
  xlabel={$x$}, ylabel={$y$},
  every axis x label/.style={at={(axis cs:4.5,0)}, anchor=west},
  every axis y label/.style={at={(axis cs:0,4.5)}, anchor=south},
  xmin=-3, xmax=5, ymin=-5, ymax=5,
  xtick={-3,-2,-1,0,1,2,3,4},
  ytick={-4,-2,0,2,4},
  tick label style={font=\footnotesize},
  width=8cm, height=8cm,
  thick,
]
  % 抛物线
  \addplot[blue, thick, domain=-1.5:3.5, samples=100] {x^2-2*x-3};
  \node[blue, right, font=\footnotesize] at (axis cs:3.2,4.0) {$y=x^2-2x-3$};

  % 固定点A和B
  \fill (axis cs:-1,0) circle(2.5pt);
  \node[above left, font=\small] at (axis cs:-1,0) {$A(-1,0)$};
  \fill (axis cs:3,0) circle(2.5pt);
  \node[above right, font=\small] at (axis cs:3,0) {$B(3,0)$};
  % AB线段
  \draw[thick] (axis cs:-1,0) -- (axis cs:3,0);

  % 情形1: 直角在A, P在x=-1竖线, AP⊥AB => P在(-1,y)
  % AP=AB=4, P(-1,-4)
  \fill[red] (axis cs:-1,-4) circle(2.5pt);
  \node[left, red, font=\small] at (axis cs:-1,-4) {$P_1(-1,-4)$};
  \draw[red, dashed] (axis cs:-1,0) -- (axis cs:-1,-4);
  \draw[red, dashed] (axis cs:-1,-4) -- (axis cs:3,0);
  % 直角标记在A
  \draw[red] (axis cs:-1,-0.4) -- (axis cs:-0.6,-0.4) -- (axis cs:-0.6,0);

  % 情形2: 直角在B, P在x=3竖线, BP=AB=4, P(3,-4)
  \fill[green!60!black] (axis cs:3,-4) circle(2.5pt);
  \node[right, green!60!black, font=\small] at (axis cs:3,-4) {$P_2(3,-4)$};
  \draw[green!60!black, dashed] (axis cs:3,0) -- (axis cs:3,-4);
  \draw[green!60!black, dashed] (axis cs:3,-4) -- (axis cs:-1,0);
  % 直角标记在B
  \draw[green!60!black] (axis cs:3,-0.4) -- (axis cs:2.6,-0.4) -- (axis cs:2.6,0);

  % 情形3: 直角在P, P在中垂线x=1上
  % 中垂线上P(1,-4)满足PA=PB，但要检验PA⊥PB
  % 这里仅示意P在中垂线上（不满足直角条件，故画灰色虚线表示"不存在"）
  \fill[gray] (axis cs:1,-4) circle(2pt);
  \node[below right, gray, font=\footnotesize] at (axis cs:1,-4) {$P_3(1,-4)$（仅等腰）};
  \draw[gray, dashed] (axis cs:1,0) -- (axis cs:1,-4);
  \node[gray, above, font=\footnotesize] at (axis cs:1,-0.2) {中垂线};

  \node[below left, font=\small] at (axis cs:0,0) {$O$};
\end{axis}
\end{tikzpicture}
\end{document}
